%=====================================================================
% 第 4 节  配分函数与共形块
% 原文：Section 4. Partition functions versus conformal blocks
%       （arXiv:1702.06499v2）
%=====================================================================
\section{配分函数与共形块}\label{partfct}

\begin{keytakeaways}
\begin{itemize}
  \item \textbf{物理图像：}AGT 对应及其推广以配分函数为核心，而 Kapustin--Witten 处理几何 Langlands 对应的进路中根本不出现配分函数。本节要在二维 sigma 模型中为二者找到统一的表述：四维半球上的路径积分给出拓扑向量空间中的一个向量，其波函数就是瞬子配分函数；而在二维一侧，同一个配分函数可以按两种不同的 Hamilton 解释读出。
  \item \textbf{数学工具：}带 $\Om$-形变的四维半球 $B^4_{\ep_1\ep_2}$ 上的拓扑扭曲、开/闭 TQFT 的态—边界态对应、扭 $\CD$-模与 KZB 方程，以及复 Fenchel--Nielsen 坐标下特征簇的可积结构。
  \item \textbf{关键结论：}配分函数 $\CZ(a;x,\tau;\ep_1,\ep_2)$ 既可以看成 $\CH_x^{\sst(2)}=\Hom(\mathsf{B}_{\rm cc},\mathsf{L}_x^{\sst(2)})$ 中元素的波函数，也可以看成 $\CH_a^{\sst(1)}=\Hom(\mathsf{B}_{\rm cc},\mathsf{L}_a^{\sst(1)})$ 中元素的波函数；这两种 Hamilton 解释由模不变性的开 TQFT 版本相联系，正是把 AGT 一侧与几何 Langlands 一侧缝合起来的关键。
  \item \textbf{关键构造：}边界条件 $\mathsf{L}^{\sst(1)}_{a}$ 实现为特征簇 $\CM_B(G)$ 中由长度函数 $a$ 生成的 Hamilton 流的轨道，再经和乐映射与 NAH 对应搬到 $\CM_H(G)$；这解释了记号中的“第一纤维化”。
\end{itemize}
\end{keytakeaways}

\bigskip
\noindent\rule{\textwidth}{0.6pt}
\bigskip

配分函数在 AGT 对应及其各种推广中处于核心地位，但在 Kapustin 与 Witten 处理几何 Langlands 对应的进路中却并不出现。为了建立一个统一的框架，下面我们讨论如何在二维 sigma 模型所提供的有效描述中定义相应的配分函数。

\subsection{四维与二维的 TQFT 框架}

利用 $\CS$ 类理论到二维的约化，下面我们要为配分函数在以 $\CM_H(G)$ 为目标流形的二维 sigma 模型中的描述提供动因。这一描述将基于又一类边界条件，记作 $\mathsf{L}^{\sst (1)}_{a}$。记号 $\mathsf{L}^{\sst (1)}_{a}$ 的选取源自它与 Hitchin \emph{第一}纤维化的联系，这一点后文将讨论。

仿照 \cite{NW}，我们考虑带 $\Om$-形变的、定义在半球 $B_{\ep_1\ep_2}^4$ 上的拓扑扭曲 $\CS$ 类理论。拓扑扭曲的 $\CS$ 类理论为 $M^3_{\ep_1\ep_2}=\pa B^4_{\ep_1\ep_2}$ 配上一个向量空间 $\CH_{\rm top}=Z(M^3_{\ep_1\ep_2})$，这里把它等同于超荷 $Q$ 的上同调，而 $Q$ 是能够在 $B_{\ep_1\ep_2}^4$ 上保持的那个超荷。可以用四维半球 $B_{\ep_1\ep_2}^4$ 上的路径积分定义一个向量 $\Psi \in\CH_{\rm top}$。向量 $\Psi$ 的波函数 $\Psi(a)$ 可以等同于配分函数 $Z(B^4_{\ep_1,\ep_2};\mathcal{B}_a)$，后者通过在 $M^3_{\ep_1\ep_2}$ 上施加适当的、由参数 $a$ 标记的 $Q$-不变边界条件 $\mathcal{B}_a$ 来定义。这类边界条件可以等同于定义 Nekrasov 配分函数时在 $\BR^4_{\ep_1\ep_2}$ 的无穷远处所加的边界条件，它们特别地把向量多重态中各个标量的零模固定为向量 $a=(a_1,\dots,a_{3g-3+n})$ 所收集的取值。边界条件 $\mathcal{B}_a$ 定义了一族边界态 $\beta_a$，从而可以把 $Z(B^4_{\ep_1,\ep_2};\mathcal{B}_a)$ 表示为交叠 $\langle\beta_a,\Psi\rangle$。

在把 $\CS$ 类理论约化到 $2d$ 拓扑 sigma 模型时，应当得到上面引入的 4d TQFT 数据的如下表示：
\begin{itemize}
\item 向量空间 $\CH_{\rm top}\simeq Z(S^3_{\ep_1,\ep_2})\ra Z(I)$。
\item 向量 $\Psi=Z(B^4_{\ep_1,\ep_2})\ra Z(T_{\ep_1,\ep_2})\in\CH_{\rm top}$，其中 $T_{\ep_1,\ep_2}$ 是挖去“上”边之后的那个开三角形，它在拓扑上等价于 $\BR_-\times I$，并通过在 $\BR_-$ 的无穷远处添加一个点而部分紧化。$T_{\ep_1,\ep_2}$ 的边界为 $\{0\}\times I$。
\item 配分函数 $Z(B^4_{\ep_1,\ep_2};\mathcal{B}_a)\ra Z(T_{\ep_1,\ep_2};\mathsf{B}_a)$ 则与三角形 $T_{\ep_1,\ep_2}$ 相联系，其上边 $\{0\}\times I$ 被赋予边界条件 $\mathsf{L}_{a}^{\sst (1)}$。$\mathsf{L}_{a}^{\sst (1)}$ 是由赋给 $M^3_{\ep_1\ep_2}$ 的边界条件 $\mathcal{B}_a$ 经约化到一维而定义的。
\end{itemize}
这意味着瞬子配分函数 $\CZ(a;x;\tau;\ep_1,\ep_2)$ 可以由三角形上 sigma 模型的配分函数表示，而该三角形的各边分别染上 $(\mathsf{B}_{\rm cc}, \mathsf{L}_a^{\sst (1)}, \mathsf{L}_x^{\sst (2)})$。

\subsection{$\CD$-模与配分函数}

量子理论中的态 $\psi$ 抽象地由 Hilbert 空间 $\CH$ 中的向量表示。要得到一个作为波函数空间 $\psi(x)$ 的具体表示，需要引入 $\CH$ 的（Hermite）对偶 $\CH^\dagger$ 中的一族元素 $\de_x$，使得 $\psi(x)$ 可以表示为 $\langle\de_x,\psi\rangle$，其中 $\langle . , .\rangle:\CH^{\dagger}\times\CH\ra \C$ 是自然的配对。

在目前的语境下，我们可以把配分函数 $\CZ=\CZ(a;x,\tau;\ep_1,\ep_2)$ 解释为元素 $\phi_x\in \CH_{x}^{\sst(2)}:=
{\rm Hom}_{\CM_H(G)}^{}(\mathsf{B}_{\rm cc}^{},\mathsf{L}_{x}^{\sst(2)})$ 的波函数 $\phi_x(a)=\langle\de_a^{\sst(1)},\phi_x^{}\rangle$，该元素由半无限长条 $\BR_-\times I$ 上的 sigma 模型路径积分产生。“边界态”$\de_a^{\sst(1)}$ 表示的正是加在 $\{0\}\times I$ 上的边界条件 $\mathsf{L}_{a}^{\sst(1)}$。

交叠 $\langle\de_a^{\sst(2)},\phi_x^{}\rangle$ 可以表示为拓扑 sigma 模型在三角形上的一个配分函数，该三角形三条边上分别赋以三种不同的边界条件 $(\mathsf{B}_{\rm cc},\mathsf{L}_{x}^{\sst(2)},\mathsf{L}_{a}^{\sst(1)})$。时间反演对称性把 $\CZ$ 的复共轭 $\CZ^*$ 等同于拓扑 sigma 模型在三角形上的一个配分函数，只是其边界条件 $(\mathsf{B}_{\rm cc}^{},\mathsf{L}_{a}^{\sst(1)},\mathsf{L}_{x}^{\sst(2)})$ 以不同的次序出现。$\CZ^*$ 具有一个 Hamilton 解释：它是元素 $\vf_a\in \CH_{x}^{\sst(1)}:=
{\rm Hom}_{\CM_H(G)}^{}(\mathsf{B}_{\rm cc}^{},\mathsf{L}_{a}^{\sst(1)})$ 的波函数 $\vf_a(x)=\langle\de_x^{\sst(2)},\vf_a^{}\rangle$，而该元素由半无限长条 $\BR_-\times I$ 上的 sigma 模型路径积分产生。

两个空间 $\CH_{a}^{\sst(1)}$ 与 $\CH_{x}^{\sst(2)}$ 都具有 $\mathcal{A}_{dR}^{\hbar}$-模结构，它来自代数 $\mathcal{A}_{dR}^{\hbar}$ 在 ${\rm Hom}_{\CM_H(G)}^{}(\mathsf{B}_{\rm cc}^{},\mathsf{B})$ 上的自然作用。上文已论证，这一模结构与由仿射代数 Ward 恒等式导出的、共形块空间上的模结构相重合。由此得到的扭 $\CD$-模结构，与手征配分函数（对应于仿射李代数的共形块）所满足的 KZB 方程所定义的结构相重合。配分函数 $\CZ$ 是 KZB 方程的解。$\CD$-模 $\CH_{a}^{\sst(1)}$ 与 $\CH_{x}^{\sst(2)}$ 可以由相应的波函数空间完全刻画，而这些波函数在物理上就表现为配分函数 $\CZ$。

在下面将要讨论的 $\ep_2\ra 0$ 极限下，我们将提出：$\CH_{{a}}^{\sst (1)}$ 出现在 Kapustin--Witten 处理几何 Langlands 对应的进路中，而 $\CH_x^{\sst(2)}$ 在 AGT 对应的各种推广的语境下似乎是更自然要考虑的空间。由模不变性的一个开 TQFT 版本所给出的、配分函数 $\CZ(x,{a},\tau;\ep_1,\ep_2)$ 之两种 Hamilton 解释之间的关系，正是这两个故事之间的关键纽带。

\subsection{膜 $\mathsf{L}_{a}^{\sst (1)}$ 的 sigma 模型描述}\label{branes1}

现在我们来给出边界条件 $\mathsf{L}_{a}^{\sst(1)}$ 作为 $\CM_H(G)$ 中 Lagrangian 子流形的一个描述。为后文之便，我们把指代用于定义膜 $\mathsf{L}_{a}^{\sst(1)}$ 的 $\CS$ 类理论之 Lagrangian 表示的标号 $\si$ 并入记号，改记为 $\mathsf{L}_{\si,a}^{\sst(1)}$。我们提出：膜 $\mathsf{L}_{\si,a}^{\sst(1)}$ 可以描述为特征簇 $\CM_B(G)$ 中由一族长度函数 $a=(a_1,\dots,a_{3g-3+n})$ 所生成的 Hamilton 流的轨道；这一族长度函数构成复 Fenchel--Nielsen 坐标系{$^{g)}$}中互相对易（Poisson 对易）的那一半，而该坐标系曾在 \cite{NRS} 的一个相关语境中被使用。这些轨道经和乐映射与 NAH 对应映到 Hitchin 模空间 $\CM_H(G)$。从特征簇 $\CM_B(G)$ 的可积结构来看，膜 $\mathsf{L}_{{\ep},a}^{\sst(1)\si}$ 是自然的对象。不同的复 Fenchel--Nielsen 坐标系由裤分解 $\si$ 标记。

我们的提议基于 {\crd \cite{DMO}} 中已知的关系：$\CS$ 类理论中的 Wilson 圈可观测量与六维 $(2,0)$ 理论中支撑在两个圆 $S^1_\CS$ 与 $S^1_C$ 之积上的余四维缺陷相关联，而这两个圆分别位于 $\CM^4$ 与 $C$ 之中。这些关系可以推广到存在面算符的情形。若把注意力集中在与 $\CM^4$ 中面算符不共享支撑的 Wilson 圈上，则可如前论证：Wilson 圈可观测量的本征值可以用标量零模 $a$ 来表示，而 $a$ 正是出现在瞬子配分函数 $\CZ({\crd a;x},\tau;\ep_1,\ep_2)$ 自变量中的量。

另一种做法是把 \cite{AGGTV} 中的论证推广到带有额外余二维面算符的情形，而这可以建立在第 \ref{Hecke1} 节所述结果的基础之上。

\bigskip
\noindent\rule{\textwidth}{0.6pt}

% [Tutor Note: 第 1 小节 itemize 第三条内部自相矛盾：先写 $Z(T_{\ep_1,\ep_2};\mathsf{B}_a)$，紧接着又说上边赋的是边界条件 $\mathsf{L}_{a}^{\sst (1)}$，并通篇只用 $\mathsf{L}_a^{\sst (1)}$。按上下文 $\mathsf{B}_a$ 应为 $\mathsf{L}_a^{\sst (1)}$ 的笔误（$\mathsf{B}$ 是膜/boundary condition 的旧记号，$\mathcal{B}_a$ 则是四维边界条件）；译文按原文保留，读者不要据此认为存在两个不同的边界条件。]
% [Tutor Note: 第二小节的指标有系统性的错位：由 $L_a^{\sst(1)}$ 定义的边界态本应记作 $\de_a^{\sst(1)}$，但正文写成交叠 $\langle\de_a^{\sst(2)},\phi_x\rangle$，上标 $(2)$ 与 $\phi_x$ 一侧重复；同样，$\vf_a$ 被声明属于 $\CH_{x}^{\sst(1)}:=\Hom(\mathsf{B}_{\rm cc},\mathsf{L}_{a}^{\sst(1)})$，下标 $x$ 应为 $a$（上标 $(1)$ 正确）。后文又出现正确的 $\CH_{a}^{\sst(1)}$，说明是排印疏漏。核对波函数的自变量与下标是否匹配（$\phi_x(a)$ 对 $\CH^{(2)}_x$、$\vf_a(x)$ 对 $\CH^{(1)}_a$）是理解本节的关键。]
% [Tutor Note: 配分函数自变量的分隔符在原文中不统一：$\CZ(a;x;\tau;\ep_1,\ep_2)$（分号分隔 $x$ 与 $\tau$）、$\CZ(a;x,\tau;\ep_1,\ep_2)$、$\CZ(x,{a},\tau;\ep_1,\ep_2)$（$x$ 与 $a$ 的次序颠倒且全用逗号）三种写法并存。按本文约定，分号用于区分四维数据 $(a;\tau,\ep)$ 与二维/共形块数据 $x$，因此第二种写法是规范的；第三种不仅是书写次序问题，其“两种 Hamilton 解释”论断恰恰要求把 $a$ 与 $x$ 视为分别属于 $\CH^{(1)}$ 与 $\CH^{(2)}$ 的两套标记，混写时须小心不要颠倒 $\langle\de_a,\phi_x\rangle$ 与 $\langle\de_x,\vf_a\rangle$ 的配对方向。]
% [Tutor Note: 膜 $\mathsf{L}_{{\ep},a}^{\sst(1)\si}$ 的记号中多出一个 $\ep$ 下标：同段前后文均只有 $\mathsf{L}_{\si,a}^{\sst(1)}$ 与 $\mathsf{L}_{a}^{\sst(1)}$，且 $\ep$ 在全文专指 $\Om$-形变参数 $\ep_1,\ep_2$，此处无任何作用。对照第 7 节对同一对象的引用（那里写作 $\mathsf{L}_{\si,a}^{\sst (1)}$）可以确认这是排印时残留的标记。]
% [Tutor Note: 4d TQFT 段落里“$\CH_{\rm top}=Z(M^3)$ 等同于 $Q$ 的上同调”这一等同需要扭曲使 $Q$ 幂零、并取 $Q$-上同调才是拓扑不变量（Witten 型论证）；而“向量 $\Psi$ 由半球上的路径积分给出、其波函数即 $\langle\beta_a,\Psi\rangle$”是开/闭 TQFT（Atiyah 公理）中态—边界态的通用构造，并不需要 $\CS$ 类理论的具体细节。真正被跳过的是末段“两种 Hamilton 解释由模不变性的开 TQFT 版本相联系”：这本质上是开弦—闭弦对偶/Cardy 条件在有边界与缺陷情形下的推广，其成立条件（边界态构成完备基、模变换的收敛性）在无穷维情形并未验证，作者把它当作出发点而非结论，引用时不宜说成已证明。]
% [Tutor Note: 记号 $S^1_\CS$ 与 $S^1_C$ 的下标在原文中并不自洽：全文的宏 \CS 定义为花体 $\mathcal{S}$（class S 的 S），故这里排出来是 $S^1_{\mathcal S}$，但从上下文（complex CS theory、Chern--Simons）看作者想写的是 Chern--Simons 的“CS”。译文按原文保留该记号。读者须按语境判断 $\CS$ 指 class S 还是 Chern--Simons：本节其余各处（$\CS$ 类理论）都是前者，只有 $S^1_\CS$ 这一处是后者，核对余四维缺陷所缠绕的两个圆时不要混淆。]
